Cross-sectional HIV incidence estimation from recency assays underpins counterfactual-placebo designs in prevention trials, where randomised placebo control is no longer ethically available. These estimators are defined on a population eligible at the time of survey, yet eligibility can change between infection and sampling --- through incarceration, displacement, hospitalisation, migration or death. Existing frameworks treat survey-time eligibility as a primitive and bound the problem by assuming that few individuals move, which is not the case in the populations for which counterfactual designs are most often required.

We retain the estimand, recency function and assumptions of the established framework and refine the eligibility indicator as the observable marginal of a finite-state process comprising an observable state, temporarily unobservable states permitting return, and an absorbing state. Allowing HIV acquisition in every living state rather than only while observable, we obtain the probability limit of the adjusted estimator as an integral of the recency function against a historical observability weight.

Temporary loss of eligibility is not inherently biasing. Under stable living-state composition, a demographically stationary observable susceptible pool, state-invariant acquisition, infection-independent movement and no absorbing loss, losses and returns cancel exactly --- for any admissible recency function and at any occupancy of unobservable states. Cancellation holds within stationary strata and therefore under arbitrary heterogeneity in movement propensity, so concentration of movement in a high-propensity minority does not generate bias. Bias requires a specific symmetry failure: absorbing loss, state-dependent acquisition, infection-dependent movement, or non-stationarity. At sourced mortality of 0.040 per year the attenuation is 2.1\%. State-dependent acquisition can attenuate or inflate, is the governing parameter, and is not identifiable from routine surveillance; we report it as a sensitivity axis rather than a corrected estimate. Analytic results were verified against an independently written generative simulator sharing no code with the derivation.

Occupancy of temporarily unobservable states cannot by itself establish that a cross-sectional incidence estimate is biased, or in which direction. Absorbing and temporary loss are not interchangeable, and correction is warranted only where a named condition fails.

\textbf{Keywords:} cross-sectional HIV incidence; recent infection testing algorithm; eligibility; selection bias; identifiability; counterfactual placebo; people who inject drugs

\begin{center}\rule{0.5\linewidth}{0.5pt}\end{center}
